\documentclass[a4paper]{article}
\usepackage[pdftex]{graphicx}
\usepackage[utf8]{inputenc}
\usepackage{enumerate}
\usepackage{siunitx}
\sisetup{locale=DE} 
\usepackage{href-ul}
\hypersetup{
	colorlinks=true,
	linkcolor=blue,
	urlcolor=blue}
\usepackage{icomma}
\usepackage{amssymb}
\usepackage{geometry}
\geometry{a4paper, top=15mm, left=15mm, right=15mm, bottom=15mm,
	headsep=10mm, footskip=12mm}

\begin{document}
{\bf Kreisbewegung und Gravitation}
\begin{enumerate}[1.]

%Aufgabe 3
\begin{minipage}{0.65\textwidth}
{\bf \item Rotierende, senkrechte Scheibe}\\
Eine Kugel befindet sich in einem offenen Fach auf einer Scheibe, die um eine waagrechte Achse rotiert. Wie groß muss ...

\begin{enumerate}[a)]
\item die Winkelgeschwindigkeit
\item die Drehfrequenz
\item die Umlaufdauer
\item die Bahngeschwindigkeit
\end{enumerate} 

... in Abhängigkeit vom Scheibenradius $R$ (Für $R=\SI{6,0}{\meter}$) mindestens sein, damit die Kugel in ihrem Fach bleibt?
\end{minipage}
\hfill
\begin{minipage}{0.35\textwidth}
	\includegraphics[width=6 cm]{scheib415}
\end{minipage}



\begin{minipage}{0.7\textwidth}
{\bf \item Halbes Looping}\\
Wie stark muss die Feder mit der Härte $D=\SI{1000}{\newton\over \meter}$ gespannt werden, damit sie dem Wagen genug Geschwindigkeit verleiht, um im Punkt A anzugelangen, ohne den Kontakt zur Wandung zu verlieren? Der Radius des Halbbogens sei $R=\SI{1,5}{\meter}$, die Masse des Wagens $\SI{1,2}{\kilogram}$
\end{minipage}
\hfill
\begin{minipage}{0.3\textwidth}
	\includegraphics[width=4 cm]{loop415}
\end{minipage}

{\bf \item Mit welchem Winkel zur Vertikalen legt sich ein Radfahrer}

\begin{enumerate}[a)]
	\item  in die Kurve mit Radius $R=\SI{30}{\meter}$, wenn er $\SI{6,0}{\meter\over\second}$ schnell fährt?
	\item maximal in eine Kurve ohne wegzurutschen, wenn der Haftreibungskoeffizient $\mu =0,70$ beträgt?
\end{enumerate}

{\bf \item Wie schnell darf Lunerald mit seinem Sportwagen} in eine Kurve mit Radius $r=\SI{50}{\meter}$ maximal fahren, wenn der Haftreibungskoeffizient $\mu=0,80$ beträgt?
	
{\bf \item Kettenkarussell} Wie groß ist die Bahngeschwindigkeit $v$ der ,,Passagiere", wenn der  Winkel $\alpha=20^\circ$ beträgt? Die  Radiuslänge sei $r=\SI{3,0}{\meter}$ und die Kettenlänge $l=\SI{3,5}{\meter}$  Welchen Wert hat dann die Umlaufdauer $T$?

\includegraphics[width=12 cm]{karu415}

\end{enumerate}

\fbox{
	\begin{minipage}{0.5\textwidth}
		Zur Lösung bitte \href{https://www.okuyakl.de/phys/p10krbL415/ll415.pdf}{hier klicken} oder den QR-Code scannen.\\
	Weitere Arbeitsblätter gibt es unter 
	
	\href{https://www.okuyakl.de}{www.okuyakl.de}
	\end{minipage}
	\hfill
	\begin{minipage}{0.4\textwidth}
		\includegraphics[width=1.5 cm]{../../viecher/zwe03}
		\includegraphics[width=3 cm]{qr415}
		\includegraphics[width=2 cm]{../../viecher/afanticon1}
		
	\end{minipage}}
\end{document}